\chapter[Construction Disk]{Making and Using a „Construction Disk“}
\label{sec:ConstructionDisk}
\newpage

\section{The construction}
Using a construction disk helps you to discover what changes you can make to improve your boomerang. With it you can better understand why your boomerang flies like it does.
You will need a wooden board or thick cardboard approx. 50 X 100 cm. Draw circles and lines onto the board, according to figure \ref{fig:CD}.
\begin{figure}[H]
\centering
\includegraphics[height=30em]{67_KONST.eps}
\caption{Construction disk}
\label{fig:CD}
\end{figure}


\section{Finding centre of rotation}
First you will have to find the boomerang’s centre of rotation. The simplest way to do this is to place the boomerang on the edge of a sharp-edged table with one tip near the edge of the table. Balance the boomerang so that it just starts to tip, then mark the boomerang directly above the edge of the table. Repeat again with the other tip of the boomerang near the edge of the table (see figure \ref{fig:CoR}). Connect the points with tape (see figure \ref{fig:CoR}). Draw lines on the tape from one arm to the other along the line on which the boomerang first balanced and then along the second balance line. The cross point represents the boomerang’s centre of rotation. 

\begin{figure}[H]
\centering
\includegraphics[height=30em]{68_SCHW.eps}
\caption{Finding the center of rotation}
\label{fig:CoR}
\end{figure}

\section{Placing boomerang on construction disk}
Now try to estimate the ratio of $V_t$ (forward speed) and $V_r$ (rotation speed). In the following example, I am assuming that $V_t = V_r$ which means that the forward speed equals the rotation speed.

Place the boomerang with the centre of rotation on the vertical line V so that one wing tip is just above the instant centre on the construction disk (see figure \ref{fig:BoomerangCD}). The lines represent the direction and relative speed of the air flow on the wing. The lines which are further away from the instant centre represent a higher speed of air flow.
If you turn the boomerang around its centre of rotation in small steps, you can see that the air flows across the wings at different angles as it rotates.

\begin{figure}[H]
\centering
\includegraphics[height=30em]{71_KOBUM.eps}
\caption{Placing boomerang on Construction disk}
\label{fig:BoomerangCD}
\end{figure}
\newpage

\section{Find the following points}
\begin{itemize}
\item maximum lift of each wing, when each is at a right angle to the air flow. 
Turn the boomerang to the left (to the right for left-handers) until the tip (the last 2 inches) of one wing lies at a right angle to the circle lines and mark the appropriate point (see figure \ref{fig:MaxLiftWing}). Do the same to all other wings.

\begin{figure}[H]
\centering
\includegraphics[width=20em]{73_10CM.eps}
\caption{Maximum lift of wing}
\label{fig:MaxLiftWing}
\end{figure}

\item maximum lift of each wing, caused by dihedral (if present).
Turn the boomerang to the left (to the right for left-handers) until the tip of one wing (the last 2 inches) lies parallel to the circle lines and mark the appropriate point (see figure \ref{fig:MaxDiLift}). Do the same to each other wing.  Note whether the dihedral is positive or negative.

\begin{figure}[p]
\centering
\includegraphics[height=20em]{72_LBIEG.eps}
\caption{Max lift caused be dihedral}
\label{fig:MaxDiLift}
\end{figure}

\item maximum lift of the elbow (if present)
Turn the boomerang so the centre of the elbow lies at a right angle to the circle lines (see figure \ref{fig:MaxElift}). This is not as easy. Mark whether the lift is positive or negative created by the elbow.

\begin{figure}[p]
\centering
\includegraphics[width=20em]{73AELLBO.eps}
\caption{Max lift on elbow}
\label{fig:MaxElift}
\end{figure}

\end{itemize}
\newpage
To determine the average point of lift, start by assuming that the wings have the same profile, width etc. Make adjustments if you think it is necessary. If arm 2 is hit by a faster air flow than arm 1, for example, arm 2 will have more lift. This will move the average maximum point of lift slightly towards arm 2. Increases or decreases of lift caused by undercutting or carving out must also be considered.

Determining the average maximum point of lift from all points on the boomerang is a difficult mathematical task. To simplify, just try to find the average point of lift resulting from the lift at points $A_1$ and $A_2$. 

Once the average maximum point of lift has been determined, mark 90° ahead in the direction of rotation as the actual point of attack. If it lies exactly on the horizontal line, the boomerang will have little or no lay-over. This would be ideal for ``Fast Catch''. If it lies below the horizontal line, the boomerang has lay-over. The lower the point the more lay-over it has and the higher the boomerang will tend to fly. If it lies above the line, it will be very difficult to get the boomerang to return, because it will roll inward and dive.

At points $L_1$, $L_2$ and $E_1$ it is very difficult to estimate the lift. Generally, positive dihedral increases lay-over and negative lift of the elbow decreases the lay-over. The value of these numbers is not important. However, you can see what changes might be expected if you change the dihedral or alter the airfoil at a certain place on one of the arms or at a place on the elbow. 

Try to analyze, for example what would happen if the lift on the middle of arm 2 was increased significantly. Analyze that point during all of its 360° degrees of rotation. Then analyze another point on the boomerang throughout its rotation. Then try placing the boomerang at a point further from the instant centre. Typically, boomerangs have significantly more forward speed than rotation speed when thrown. The instant centre is therefore at least one boomerang span from the centre of rotation. At the end of their flight, $V_t$ approaches 0. This is why lay-over is different early in the flight than it is late in the flight.

By experimenting, with the construction disk as your guide, you should be able to find points on the boomerang that you can carve that will have significant effect during the first part of the flight, without affecting the latter part of the flight and vice versa.



